\chapter{Teaching a route}
\label{ch:u-routes}

A route is a walk the robot repeats. You teach it once, by walking the robot along it, and mark
the stop points where it will stand still and inspect.

\section{Where routes go}

\begin{itemize}
  \item Along walkways and roads people may use while the equipment is energized, outside every
    minimum approach distance.
  \item On even, solid ground the robot can see: no steps down, no holes, no loose gravel slopes.
  \item With room: at least 1~m wider than the robot (0.47~m) where it passes things, so it can
    step around an obstacle.
  \item Starting at the dock, so every patrol begins where the robot was started.
\end{itemize}

\section{Teach}

\begin{enumerate}
  \item REMOTE (key \key{A}). The robot stands on the dock, \ui{Start pose} in the top bar.
  \item \ui{Record}: \ui{Recording a route}.
  \item Walk the route slowly and smoothly, as you want it repeated. The robot remembers its path
    every few centimetres.
  \item At each place to inspect, stop the robot where the cameras see the equipment and press
    \ui{Point}: \ui{Stop point, inspect 10 s}. The robot will stand there on every patrol.
  \item At the end, \ui{Finish}: \ui{Route saved}. Its name is the date and the time.
\end{enumerate}

The route is saved on the robot and sent to NETRA by itself (\menupath{Routes}), now or as soon as
NETRA is reachable.

\begin{tipbox}[Stop points]
Stop where the robot sees the equipment straight and close enough, without leaning into any
approach distance. The inspection points are then taught for that stop
(Chapter~\ref{ch:u-inspection}); a stop with none simply waits its time.
\end{tipbox}

\section{Tags}

Where the site has markers (tags) on walls or posts, put the robot in view of one near the start
and press \ui{Tags}: the robot records where it is. On later patrols, each time it sees that tag
it corrects the small drift of its position. Ask the technician where tags help.
